\section{Limitations}
\label{sec:limitations}

These are ordered by how much they could move a conclusion.

\textbf{1. The memorisation check could not be validated.} Protection against a
model recalling an outcome rests on the release-date rule and the memory probe.
A model released long after \winA settled (\mbox{GLM-5.3 Flash}, 26 August
2026) was flagged at 1.4 per cent on 500 \winA rows, less than the 1.6 to 3.4
per cent for the eligible models on the same rows. The probe is therefore a
filter of unknown power, and the leakage argument rests on release dates alone
(Appendix~\ref{app:memory}).

\textbf{2. Short, sports-heavy windows.} Sports make up 82, 73 and 81 per cent
of primary rows, and our findings may not transfer to slower political or
macroeconomic questions. The gap is larger on non-sports rows
(Appendix~\ref{app:nonsports}), so the sports concentration is not what produces
it; but the non-sports evidence is thin.

\textbf{3. One prompt, one sample, and what that means.} Every model answered one
prompt once, at temperature 0 where the host accepted it. We did not vary the prompt
wording, so the measured agreement may overstate what would
remain under varied prompts: we can say the models move together given a common
question but cannot separate shared judgment from shared phrasing. The
same-developer comparison does not settle it either way, for the reasons
Section~\ref{sec:results-correlated} gives. Nor can we say how much of each
model's noise is specific to the one wording and draw it got; we did not average over
prompt variants or repeated samples.

\textbf{4. Thinking budgets are not comparable.} At the 1{,}000-token cap sit
99 per cent of \mKimiKTwoFive's and \mKimiKTwoSix's rows, 86 to 92 per cent of
\mDeepSeekVThreeTwo's and 46 to 50 per cent of \mDeepSeekVFourPro's (its median
over the whole run, primary and non-primary rows together, sits at the cap in
Table~\ref{tab:roster}), so we cannot say
what they would do uncapped. The model-versus-market comparison does not depend
on cross-model comparability, but a capped model might do better uncapped.

\textbf{5. Hosts and numeric precision differ.} Open-weight models are served at
whatever precision the host offers, including int4 for \mKimiKTwoFive, and one
host truncated \mDeepSeekVFourPro's output at 3{,}000 tokens. A differently
served copy might score differently.

\textbf{6. Question text on \winB and \winC was captured after settlement.}
Exchanges can edit text after listing, so equality with what a trader saw is
assumed, not verified. \winA's text is verified, being captured live. Had an
exchange clarified a rule mid-life, our models would have seen the clarified
version and the crowd would not, which favours the models.

\textbf{7. The news arm rests on few rows.} Only 1{,}434, 563 and 752 primary
rows carry news, and the analyses run on the test-split share of those.

\textbf{8. The price is not a certified probability, and the midpoint is not
executable.} The accuracy comparison is against the midpoint of the best bid and
ask, a standard summary but not a price anyone can trade at; the trading
simulation fills at the quoted bid and ask with fees charged, and the two should
not be read as the same exercise. The price bounds the mean belief only
within an interval, wide except near even money~\citep{manski2006}, though in
practice prices are usually close to mean beliefs~\citep{wolfers2006}. Our
claims are about the price as a forecast; a reader who wants ``collective human
belief beats the models'' needs the extra assumption.

\textbf{9. The prompt is price-blind.} Our results say nothing about how a model
would trade with the price in view.

\textbf{10. Trading is simulated from a snapshot.} There is no queue position,
price impact, partial fill or later opportunity to trade, so returns are an
optimistic bound. This matters most for the against-the-crowd rows where the
losses are largest.

\textbf{11. \winC cannot be fully reproduced from the release.} \kalshi's terms do
not permit redistributing market data, so \winC ships as identifiers and model
outputs only.

\textbf{12. No non-LLM sports baseline.} We have not run a simple non-LLM
baseline on the sports rows, which would bound how much of the market's advantage
needs no language model.
